\section{Limitations and Future Work}
\label{sec:limitations}

\subsection{Limitations}

\paragraph{Statistical power and panel size.}
The per-region analysis rests on four anchor datasets, and energy's spacing is partly
degenerate (discrete features collapse the deciles). The dataset-level correlations use
ten to twelve datasets with $p$-values uncorrected over many metric pairings, so some
associations partly reflect dataset size or dimensionality rather than spacing per se. The
strong, robust coupling is specifically \texttt{force\_ii\_tmag}; the other modes are
directionally consistent but weaker at this sample size. The teacher-edge correlations use
only the four anchor datasets and are not statistically significant. The ``less-is-more''
phenomenon is robust on just a few heterogeneous datasets (\texttt{nyse\_log},
\texttt{realestate}, \texttt{metro}); airquality's instance is seed-fragile and should not
be treated as canonical.

\paragraph{Protocol and configuration scope.}
All results are obtained within one masked-distillation protocol---nested structured
$k$-NN holes with a fixed teacher (80\% kept) and student (70\% kept)---so other scarcity
structures, mask fractions, or nesting depths may behave differently. The teacher is a
single frozen model rather than an ensemble, and the core hyperparameters ($T$, clause
count, specificity, blend factor $f$, teacher-firing probability $p_{\text{teacher}}$, and
epoch budget) are held fixed rather than swept, so the reported effect sizes are specific
to this operating point.

\paragraph{Mechanism and measurement scope.}
The mechanisms act on the per-sample error sign and magnitude; we did not explore
scheduling $p_{\text{teacher}}$ over training, per-region teacher trust, or combinations of
modes. Our spacing metrics are computed in input (feature) space, and \texttt{spc\_err\_strength}
is a derived heuristic rather than a first-principles quantity. Finally, while the
clause-level feedback makes the mechanism inspectable, our interpretability claims are
qualitative: we did not quantify which clauses or literals the teacher actually changes.

\subsection{Future Work}

Several directions follow directly from these limitations.

\begin{itemize}
  \item \textbf{Spacing definitions.} Separate input-space spacing from joint
    input--target $(x,y)$ spacing to pin down the operative variable, and compare
    alternative density and heterogeneity measures against \texttt{spc\_err\_strength}.
  \item \textbf{Broader, more rigorous evaluation.} Extend the panel and apply
    multiplicity-corrected statistics to test whether the density--benefit law holds
    across all modes, not only \texttt{force\_ii\_tmag}.
  \item \textbf{Scarcity structures.} Test random, adversarial, class-imbalanced, and
    out-of-support (extrapolation) holes, and vary the mask fraction and nesting depth, to
    map how the global-reshaping-plus-local-repair account depends on how data is removed.
  \item \textbf{Teacher design and gating.} Explore teacher ensembles, self-distillation,
    and iterative teacher--student rounds; and, since accuracy does not predict benefit,
    build a data-geometry gate---for example a spacing-based go/no-go rule---that decides
    when to distil at all.
  \item \textbf{Mechanism extensions.} Study scheduled or learned $p_{\text{teacher}}$,
    per-region teacher trust, and combinations of teacher-bounded magnitude with adaptive
    triggering, and port the mechanisms to the weighted and convolutional variants and to
    classification Tsetlin Machines.
  \item \textbf{Interpretability and theory.} Quantify what the teacher transfers at the
    clause level---turning the ``inspectable'' claim into measurements---and develop a
    theoretical account of why global reshaping dominates and when localized in-hole repair
    emerges from the Type~I/Type~II dynamics.
\end{itemize}
